\section{Solution to The Issue of Correlated Prescales}
\label{sec:app-prw-math}
%We use boolean variables to represent the ``prescale enaction'' for a given trigger combination. %
\subsection{Preliminary Definitions, Notation, and Algebraic Identities}
To make the discussion unambiguous, we introduce (or recall) a few definitions and some notation:\\\\
\textbf{Trigger Element}: An L1 item or the HLT part of a trigger chain. Notice that the prescales are applied to an L1 item or the HLT part of a trigger chain, not to a trigger chain as a whole. The prescale of a trigger chain is thus emergent from the prescales of its trigger elements. %
We will use a \texttt{mathcal} uppercase letter (e.g., $\mathcal{X}$) to denote a generic trigger chain or a logical composition of generic trigger chains (referred to as a trigger combination hereafter) and a standard uppercase letter (e.g., $X$) to denote a trigger element. %
\\\\
\textbf{Trigger Decision}: This is simply the trigger decision for a given trigger combination \textbf{before} applying prescales. We will use $t^i_\mathcal{X}$ to denote the trigger decision for a trigger combination $\mathcal{X}$ in an event $i$.\\\\
\textbf{Prescale Decision}: If a trigger combination is precluded from firing due to the application of prescales, then the prescale decision for that trigger combination is $0$; otherwise $1$. We will use $u_\mathcal{X}^i$ to denote the prescale decision for a trigger combination $\mathcal{X}$ in an event $i$.
\\\\
\textbf{Passing Set of a Trigger Combination}: The set of all events for which the trigger decision is $1$ for a given trigger combination. We will use $T_\mathcal{X}$ to denote the passing set of a trigger combination $\mathcal{X}$, i.e., $T_{\mathcal{X}}=\{i\ \vert\ t^i_\mathcal{X}=1\}$. \\\\
\textbf{Unprescaled Set of a Trigger Element}: The set of all events for which the prescale decision is $1$ for a given trigger element. We will use $U_X$ to denote the unprescaled set of a trigger element $X$, i.e., $U_X=\{i\ \vert \ u^i_X=1\}$.
\\\\
We adopt the notation wherein the logical AND between triggers is expressed as a product, and the logical OR between triggers is expressed as a sum:
\begin{subequations}
\label{eq:composition-set-operations}
\begin{align}
    U_{\mathcal{X}+\mathcal{Y}}&\equiv\{i\ \vert\ u^i_\mathcal{X}\lor u^i_{\mathcal{Y}}=1\}=U_\mathcal{X}\cup U_\mathcal{Y}\\
    U_{\mathcal{XY}}&\equiv\{i\ \vert\ u^i_\mathcal{X}\land u^i_{\mathcal{Y}}=1\}=U_\mathcal{X}\cap U_\mathcal{Y}\\
    T_{\mathcal{X}+\mathcal{Y}}&\equiv\{i\ \vert\ t^i_\mathcal{X}\lor t^i_{\mathcal{Y}}=1\}=T_\mathcal{X}\cup T_\mathcal{Y}\\
    T_{\mathcal{XY}}&\equiv\{i\ \vert\ t^i_\mathcal{X}\land t^i_{\mathcal{Y}}=1\}=T_\mathcal{X}\cap T_\mathcal{Y}
\end{align}
\end{subequations}
Or, in general, if a trigger combination $\mathcal{X}$ can be expressed as a logical composition $f$ of an ordered set of trigger chains $[\mathcal{A}_k]$, i.e., $\mathcal{X}=f([\mathcal{A}_k])$, then the same logical composition also describes $t_\mathcal{X}^i$, $u_\mathcal{X}^i$ in terms of $[t_{\mathcal{A}_k}^i]$, $[u_{\mathcal{A}_k}^i]$ respectively. %
For example, if $\mathcal{A}=AX+BY$ then $t^i_\mathcal{A}=(t^i_A\land t^i_X)\lor (t^i_B\land t^i_Y)$, etc. %
This can be made more explicit: a generic trigger combination that can be created out of the $n$ trigger chains in the ordered set $[\mathcal{A}_k]$ (where $k\in \mathbb{N}_{\leq n}$) can be expressed as
\begin{align}
\mathcal{X}&=\sum_{l=1}^{2^n} e_{l}\bigg(\Pi_{k=1}^{n}\mathcal{A}_k^{c_{kl}}\bigg)
\end{align}
where $c_{kl}\in\{0,1\}$ and $e_l\in\{0,1\}$, i.e., the relevant logical composition is parameterized by $2^n+n2^n$ boolean parameters. To say that the same logical composition that describes $\mathcal{X}$ in terms of $[\mathcal{A}_k]$ also describes $t_\mathcal{X}^i$, $u_\mathcal{X}^i$ in terms of $[t_{\mathcal{A}_k}^i]$, $[u_{\mathcal{A}_k}^i]$ respectively is to say that 
\begin{subequations}
\label{eq:generic-composition-boolean-operations}
\begin{align}
    t_\mathcal{X}^i&=\lor_{l=1}^{2^{n}} e_l\bigg(\land_{k=1}^{n}{(t_{\mathcal{A}_k}^i)}^{c_{kl}}\bigg)\\
    u_\mathcal{X}^i&=\lor_{l=1}^{2^{n}} e_l\bigg(\land_{k=1}^{n}{(u_{\mathcal{A}_k}^i)}^{c_{kl}}\bigg)
\end{align}
\end{subequations}
We now introduce a few more definitions:\\\\
\textbf{Unprescaled Passing Set of a Trigger Element}: For a trigger element $X$, this is the set of all events where both the prescale decision and the trigger decision are $1$. We will use $V_X$ to denote the unprescaled passing set of a trigger element $X$, i.e., $V_X=\{i\ \vert\ u^i_X\land t^i_X=1\}$.\\\\
\textbf{Unprescaled Passing Set of a Trigger Combination}: For a trigger combination $\mathcal{X}=f([A_k])$ where $A_k$s are trigger elements and $f$ is some logical composition thereof, the unprescaled passing set of $\mathcal{X}$ is given by $V_\mathcal{X}=\{i\vert f([u_{A_k}^i\land t^i_{A_k}])=1\}$. %
For a trigger element $X$, 
        \begin{align}
            V_X=U_X\cap T_X
        \end{align}
However, for a trigger combination $\mathcal{X}=f([A_k])$, it is \textbf{NOT} generally true that $V_\mathcal{X}=U_\mathcal{X}\cap T_\mathcal{X}$. %
It does hold tho that for a trigger combination $\mathcal{X}=f([A_k])$,
        \begin{align}
        V_\mathcal{X}=f([V_{A_k}])=f([U_{A_k}\cap T_{A_k}])
        \end{align}
For example, if $\mathcal{X}=A+B$ then $V_\mathcal{X}=(U_A\cap T_A)\cup (U_B\cap T_B)\neq (U_A\cup U_B)\cap (T_A\cup T_B)$. %
Finally, just like \eqref{eq:composition-set-operations}, it holds that 
        \begin{subequations}
        \label{eq:composition-set-operations-on-v}
        \begin{align}
            V_\mathcal{X+Y}&=V_\mathcal{X}\cup V_\mathcal{Y}\\
            V_\mathcal{XY}&=V_\mathcal{X}\cap V_\mathcal{Y}
        \end{align}
        \end{subequations}
\\
\textbf{Decision Sets of a Trigger Menu}: For a trigger menu constituted by an ordered set $[A_k]$ where $A_k$s are $n$ trigger elements, a trigger decision set $\mathcal{D}^{l}$ is defined as 
    \begin{align}
        \mathcal{D}^{l}\equiv\{i\ \vert\ t^i_{A_k}=d_{kl},\ \forall k\ \in \mathbb{N}_{\leq n}\} 
    \end{align}
where the parameter $l$, which indexes the set of all trigger decision sets of the menu, runs from $1$ to $2^n$ and $d_{kl}\in \{0,1\}$ (where $k\in \mathbf{N}_{\leq n}$) is the array of boolean numbers defining the decision set $\mathcal{D}^l$.  %
Notice that we can write 
        \begin{align}
            \mathcal{D}^l=\cap_{q,\tilde{q}}T_{A_q}T^c_{A_{\tilde{q}}}
        \end{align}
where $q$ runs over the solutions of $d_{kl}=1$, i.e., $q\in \{k \vert d_{kl}=1, k\in \mathbf{N}_{\leq n}\}$ and $\tilde{q}$ runs over the solutions of $d_{kl}=0$, i.e., $\tilde{q}\in \{k \vert d_{kl}=0, k\in \mathbf{N}_{\leq n}\}$.
Using this notation, we can write
        \begin{subequations}
        \label{eq:intersect-dec-and-v}
        \begin{align}             
        V_\mathcal{X}\cap \mathcal{D}^l&=f([U_{A_k}\cap T_{A_k}])\cap_{q,\tilde{q}}T_{A_q}T^c_{A_{\tilde{q}}}\\
        &=\Big(\cup_{j}e_j\big(\cap_k (U_{A_k}\cap T_{A_k})^{c_{kj}}\big)\Big)\cap_{q,\tilde{q}}T_{A_q}T^c_{A_{\tilde{q}}}\\
        &=\cup_{j}e_j\Big(\cap_k (U_{A_k}\cap T_{A_k})^{c_{kj}}\cap_{q,\tilde{q}}T_{A_q}T^c_{A_{\tilde{q}}}\Big)\\
        &= \cup_{j\in\{m\vert d_{ml}=1\}} e_j \big(\cap_k U_{A_k}^{c_{kj}}\big) \cap_{q,\tilde{q}}T_{A_q}T^c_{A_{\tilde q}}\\
        &= \cup_{j\in\{m\vert d_{ml}=1\}} e_j\big(\cap_k U_{A_k}^{c_{kj}}\big) \cap \mathcal{D}^l\\
        &= \cup_{j}d_{jl}e_j\big(\cap_k U_{A_k}^{c_{kj}}\big) \cap \mathcal{D}^l
        \end{align}
        \end{subequations} 
\\
\textbf{Reduced Unprescaled Set of a Trigger Combination}: Notice that $U_\mathcal{X}=f([U_{A_k}])=\cup_j e_j\big(\cap_k U_{A_k}^{c_kj}\big)$. Taking a cue from \eqref{eq:intersect-dec-and-v}, we define a reduced unprescaled set of a trigger combination $\mathcal{X}$ w.r.t. a given decision set $\mathcal{D}^l$ as 
    \begin{align}
        \slashed{U}_\mathcal{X}(\mathcal{D}^l)\equiv \cup_{j}d_{jl}e_j\big(\cap_k U_{A_k}^{c_{kj}}\big)
    \end{align}
Notice that almost by construction:
    \begin{align}
        V_{\mathcal{X}}\cap \mathcal{D}^l = \slashed{U}_{\mathcal{X}}(\mathcal{D}^l)\cap \mathcal{D}^l
    \end{align}
Finally, just like \eqref{eq:composition-set-operations-on-v} and \eqref{eq:composition-set-operations}, it follows that 
        \begin{subequations}
        \label{eq:composition-set-operations-on-slashed-u}
        \begin{align}
            \slashed{U}_\mathcal{X+Y}(\mathcal{D}^l)&=\slashed{U}_\mathcal{X}(\mathcal{D}^l)\cup \slashed{U}_\mathcal{Y}(\mathcal{D}^l)\\
            \slashed{U}_\mathcal{XY}(\mathcal{D}^l)&=\slashed{U}_\mathcal{X}(\mathcal{D}^l)\cap \slashed{U}_\mathcal{Y}(\mathcal{D}^l)
        \end{align}
        \end{subequations}
\subsection{Definition of Prescale Weights}
With this machinery of definitions, we are finally able to define the prescale weights. We define the \emph{\textbf{prescale weight of an event $i$ for a given trigger combination $\mathcal{X}$ as the conditional probability that it belongs to the unprescaled passing set $V_\mathcal{X}$ of the said trigger combination given the decision set $\mathcal{D}^l$ to which the event $i$ belongs}}. Before moving on, it is important to emphasize that the prescale weight of an event \textbf{does not} correspond to the probability that an event belongs to the unprescaled set of a trigger combination. For example, imagine that we use $A\vert\vert B$ to collect data with $A$ prescaled at 2 and $B$ running unprescaled. If $B$ is a ridiculously restrictive trigger and never fires, then it's clear that the prescale weight should be $1/2$ for every event that passes $A\vert\vert B$. However, every event would always belong to the unprescaled set of $A\vert\vert B$ since $B$ is running unprescaled. \\\\
To introduce one final bit of notation, we notice that a given event $i$ unambiguously belongs to some specific decision set. Thus, we might as well specify the decision set to which a given event $i$ belongs by $\mathcal{D}^i$. This notation wouldn't be useful if one needs to run over all decision sets since two different events $i,j$ might belong to the same decision set. However, as far as the concern is unambiguously specifying a decision set, this notation is good enough. Equipped with this notation, we compute the prescale weight of an event $i$ for a given trigger combination $\mathcal{X}$, $w_{\mathcal{X}}^i$ as 
\begin{align}
\label{eq:prescale-weight-def}
    w_{\mathcal{X}}^i&=\rho(V_\mathcal{X} \vert \mathcal{D}^{i})\\
    &=\frac{\rho(V_\mathcal{X}\cap \mathcal{D}^{i})}{\rho(\mathcal{D}^{i})}\\
    &=\frac{\rho\big(\slashed{U}_{\mathcal{X}}(\mathcal{D}^i)\cap \mathcal{D}^i\big)}{\rho(\mathcal{D}^{i})}\\
    &=\rho\big(\slashed{U}_{\mathcal{X}}(\mathcal{D}^i)\big)
\end{align}
where we use $\rho$ instead of $P$ to denote probabilities as we reserve the latter to denote prescales. %
We have thus shown that \textbf{the prescale weight of a trigger combination $\mathcal{X}$ in an event $i$ is simply the probability that the event $i$ belongs to the reduced unprescaled passing set of the trigger combination $\mathcal{X}$ w.r.t. the decision set $\mathcal{D}^i$ to which the event $i$ belongs.} %
Finally, for the sake of completeness, we define the effective prescale of a trigger combination in a given event as the reciprocal of the prescale weight. 
\subsection{Evaluation of the Prescale Weights}
For a trigger element $X$, we can actually evaluate the expression \eqref{eq:prescale-weight-def}:
\begin{align}
\label{eq:weight-for-trigger-elements}
    w_X^i&=\rho(\slashed{U}_X(\mathcal{D}^i))=\rho(t_X^iU_X)=t^i_X\rho(U_X)=\frac{t^i_X}{P_X}
\end{align}
where $t^i_X$ is the trigger decision for the trigger element $X$ and $P_X$ is the prescale of the trigger element $X$. For a generic trigger combination, we first want to rewrite its weight in terms of the weights of relevant trigger elements and then utilize \eqref{eq:weight-for-trigger-elements}. 
\\\\
To this end, we derive the basic rules of the composition of prescale weights. %
\subsubsection*{Rule $1$: $w^i_{\mathcal{A}+\mathcal{B}}=w^i_{\mathcal{A}}+w^i_{\mathcal{B}}-w^i_{\mathcal{AB}}$}
Utilizing the composition rules \eqref{eq:composition-set-operations-on-slashed-u} for the reduced unprescaled sets, we can write 
\begin{subequations}
\label{eq:sum-general}
\begin{align}
    w^i_{\mathcal{A}+\mathcal{B}}&=\rho_i(\slashed{U}_{\mathcal{A}+\mathcal{B}}(\mathcal{D}^i))\\
    &=\rho_i\big(\slashed{U}_{\mathcal{A}}(\mathcal{D}^i)\cup \slashed{U}_{\mathcal{B}}(\mathcal{D}^i)\big)\\
    &=\rho_i\big(\slashed{U}_{\mathcal{A}}(\mathcal{D}^i)\big)+\rho_i\big(\slashed{U}_{\mathcal{B}}(\mathcal{D}^i)\big)-\rho_i\big(\slashed{U}_{\mathcal{A}}(\mathcal{D}^i)\cap \slashed{U}_{\mathcal{B}}(\mathcal{D}^i)\big)\\
    &=\rho_i\big(\slashed{U}_{\mathcal{A}}(\mathcal{D}^i)\big)+\rho_i\big(\slashed{U}_{\mathcal{B}}(\mathcal{D}^i)\big)-\rho_i\big(\slashed{U}_{\mathcal{A}\mathcal{B}}(\mathcal{D}^i)\big)\\
    &=w^i_{\mathcal{A}}+w^i_{\mathcal{B}}-w^i_{\mathcal{AB}}
\end{align}
\end{subequations}
\subsubsection*{Rule $2$: Iff the enactment of prescales is independent between $\mathcal{A}$ and $\mathcal{B}$, $w^i_{\mathcal{A}\mathcal{B}}=w(\mathcal{A})w(\mathcal{B})$.}
Again, utilizing the composition rules \eqref{eq:composition-set-operations-on-slashed-u} for the reduced unprescaled sets, we can write
\begin{align}
\label{eq:product-independent}
    w^i_{\mathcal{A}\mathcal{B}}&=\rho_i(\slashed{U}_\mathcal{AB}(\mathcal{D}^i))\\
    &=\rho_i\big(\slashed{U}_\mathcal{A}(\mathcal{D}^i)\cap \slashed{U}_\mathcal{B}(\mathcal{D}^i)\big)\\
    &=\rho_i\big(\slashed{U}_\mathcal{A}(\mathcal{D}^i)\big)\rho_i\big(\slashed{U}_\mathcal{B}(\mathcal{D}^i)\big)\\
    &=w^i_{\mathcal{A}}w^i_{\mathcal{B}}
\end{align}
where we assume that the prescale enactment is independent between $\mathcal{A}$ and $\mathcal{B}$ in order to write the penultimate equality. %
In particular, for trigger elements $A, B$, it always holds that $\rho(\slashed{U}_A(\mathcal{D}^i)\cap \slashed{U}_B(\mathcal{D}^i))=\rho(\slashed{U}_A(\mathcal{D}^i))\rho(U_B(\mathcal{D}^i))$ unless $A, B$ belong to a coherently prescaled group\footnote{For triggers belonging to a common coherent prescale group, the enactment of prescales is fully correlated, i.e., every event either belongs to the unprescaled set of all of the triggers in the coherent prescale group or it doesn't belong to the unprescaled set of any of the triggers in the coherent prescale group.}\footnote{Unless otherwise stated, we assume that none of the triggers being discussed are in a coherent prescale group. This assumption is always true for L1 triggers. While HLT items can be in a coherent prescale group, neither of the triggers used for this analysis was in a coherent prescale group during Run~2.}. %
However, even if $A, B$ are not coherently prescaled, for the HLT chains such as $\mathcal{A}=AX$ and $\mathcal{B}=BX$, it \textbf{does not} hold that $\rho(\slashed{U}_{\mathcal{A}\mathcal{B}}(\mathcal{D}^i))=\rho(\slashed{U}_{\mathcal{A}}(\mathcal{D}^i))\rho(\slashed{U}_{\mathcal{B}}(\mathcal{D}^i))$. %
Using the relations \ref{eq:sum-general} and \ref{eq:product-independent}, the weights for any trigger combination can be computed recursively. %
For example, if we are interested in the trigger combination $\mathcal{X}=\mathcal{A}+\mathcal{B}$, we can write
\begin{align}
w^i_\mathcal{X}&=w^i_{\mathcal{A}+\mathcal{B}}\\
&=w^i_{\mathcal{A}}+w^i_{\mathcal{B}}-w^i_{\mathcal{AB}}\\
&=w^i_{AX}+w^i_{BX}+w^i_{AXBX}\\
&=w^i_{AX}+w^i_{BX}+w^i_{ABX}\\
&=w^i_A w^i_X+w^i_B w^i_X+w^i_A w^i_B w^i_X\\
&=\frac{t^i_At^i_X}{P_A P_X}+\frac{t^i_Bt^i_X}{P_B P_X}+\frac{t^i_At^i_Bt^i_X}{P_A P_B P_X}
\end{align}
Expectedly, this is different from the weight that we would have obtained if we had ignored the correlations between the prescales of $\mathcal{A}$ and $\mathcal{B}$ (arising out of the fact that they share a common L1 item), e.g., the formula utilized by the standard prescale reweighting tool \cite{Buttinger:2014726} would have given us
\begin{align}
    w^i_{\mathcal{X}}&=w^i_{\mathcal{A}}+w^i_{\mathcal{B}}-w^i_{\mathcal{AB}}\\
    &=w^i_{\mathcal{A}}+w^i_{\mathcal{B}}-w^i_{\mathcal{A}}w^i_{\mathcal{B}}\\
    &=w^i_{AX}+w^i_{BX}-w^i_{AX}w^i_{BX}\\
    &=\frac{t^i_At^i_X}{P_A P_X}+\frac{t^i_Bt^i_X}{P_B P_X}-\frac{t^i_At^i_Bt^i_X}{P_A P_B P^2_X}
\end{align}
%A more complicated example:
%For example, if the trigger combination we are interested in is $\mathcal{X}=A(X_1+X_2)+B(X_1+X_3)$, we can write $\mathcal{X}=AX_1+AX_2+BX_1+BX_3\implies$
%\begin{align}
%w(\mathcal{X})&=w(AX_1)+w(AX_2)+w(BX_1)+w(BX_3)\\
%&-w(AX_1AX_2)-w(AX_1BX_1)-w(AX_1BX_3)\\
%&-w(AX_2BX_1)-w(AX_2BX_3)-w(BX_1BX_3)\\
%&+w(AX_1AX_2BX_1)+w(AX_1AX_2BX_3)\\
%&+w(AX_2BX_1BX_3)-w(AX_1AX_2BX_1BX_3)\\
%&=w(AX_1)+w(AX_2)+w(BX_1)+w(BX_3)\\
%&-w(AX_1X_2)-w(ABX_1)-w(ABX_1X_3)\\
%&-w(ABX_1X_2)-w(ABX_2X_3)-w(BX_1X_3)\\
%&+w(ABX_1X_2)+w(ABX_1X_2X_3)\\
%&+w(ABX_1X_2X_3)-w(ABX_1X_2X_3)\\
%&=\frac{t_{A}}{P_AP_{X_1}}+\frac
%\end{align}
\subsection{Evaluation of the Prescale Weights for the Unseeded \texttt{BeeX} Triggers}
In the electron channel of the analysis, we are interested in the following trigger combination:
\begin{align}
    \mathcal{T}=A\mathcal{X}+B\mathcal{X}
\end{align}
where \begin{align}
    A&=\texttt{HLT\_e5\_lhvloose\_nod0\_bBeexM6000t}\\
    B&=\texttt{HLT\_2e5\_lhvloose\_nod0\_bBeexM6000t}\\
    \mathcal{X}&=\sum_i X_i
\end{align}
where $i$ indexes all the L1 items in the trigger menu. We can thus write 
\begin{align}
    w^i_{\mathcal{T}}&=w^i_{A\mathcal{X}+B\mathcal{X}}\\
    &=w^i_{A\mathcal{X}}+w^i_{B\mathcal{X}}-w^i_{AB\mathcal{X}}\\
    &=w^i_{A}w^i_{\mathcal{X}}+w^i_{B}w^i_{\mathcal{X}}-w^i_Aw^i_Bw^i_{\mathcal{X}}\\
    &=\big(w^i_{A}+w^i_{B}-w^i_Aw^i_B\big)\cdot w^i_{\mathcal{X}}\\
    &=\bigg(\frac{t^i_A}{P_A}+\frac{t^i_B}{P_B}-\frac{t^i_At^i_B}{P_AP_B}\bigg)\cdot w^i_{\mathcal{X}}
\end{align}
where the $w^i_{\mathcal{X}}$ term needs to be further evaluated according to the rules derived in \eqref{eq:sum-general} and \eqref{eq:product-independent}. %
However, this computation is conceptually straightforward since the enactment of prescales is always independent among all the L1 items. %
